% Data and update signals; no gradient through discrete sampling.
\begin{tikzpicture}[
  x=1mm,y=1mm,font=\figurefont\fontsize{8}{10}\selectfont,
  box/.style={draw=BookRule,fill=BookPaper,text width=26mm,
    minimum height=11mm,align=center,inner sep=2mm,line width=.7pt},
  flow/.style={-{Stealth[length=2mm]},draw=BookMuted,line width=.8pt},
  lab/.style={fill=white,inner sep=.7mm,font=\figurefont\fontsize{7}{9}\selectfont}
]
  \node[box,fill=FigureInput!10] (input) at (16,32) {输入$q$};
  \node[box,fill=FigureModel!10] (policy) at (61,32) {旧策略$\mu=\pi^{(k)}$\\逐token采样};
  \node[box] (response) at (107,32) {回答$y$、前缀$s_t$\\保留采样概率$\mu$};
  \node[box,fill=FigureLoss!10] (reward) at (145,32) {奖励模型\\$\widehat r_{\bm\omega}(q,y)$};
  \node[box] (pairs) at (145,58) {同输入偏好对\\$(y_{\mathrm w},y_{\mathrm l})$};
  \node[box] (value) at (107,0) {价值模型\\前缀基线$V_{\bm\psi}(s_t)$};
  \node[box,fill=FigureLoss!10] (token) at (145,0) {逐token奖励\\$\widetilde r_t$};
  \node[box] (adv) at (145,-29) {回报与优势\\$\widehat A_t$};
  \node[box] (update) at (61,-29) {PPO式更新\\策略；另拟合价值};
  \node[box] (reference) at (107,-58) {固定参考策略概率\\$\pi_{\mathrm{ref}}(a_t\mid s_t)$};
  \draw[flow] (input) -- (policy);
  \draw[flow] (policy) -- node[lab,above] {生成} (response);
  \draw[flow] (response) -- (reward);
  \draw[flow] (pairs) -- node[lab,right] {偏好拟合} (reward);
  \draw[draw=BookMuted,line width=.8pt] (response.south) -- (107,15);
  \draw[flow] (107,15) -- node[lab,left] {前缀} (value.north);
  \draw[flow] (107,15) -- node[lab,above] {动作与$\mu$}
    (125,15) -- (125,0) -- (token.west);
  \draw[flow] (reward) -- node[lab,right] {序列分数} (token);
  \draw[flow] (token) -- (adv);
  \draw[flow] (value.south) -- (107,-14) -- (125,-14)
    -- (125,-29) -- (adv.west);
  \draw[flow] (adv.south) -- (145,-44)
    -- node[lab,above] {局部更新信号} (61,-44) -- (update.south);
  \draw[flow] (reference.east) -- node[lab,above] {参考概率}
    (167,-58) -- (167,0) -- (token.east);
  \draw[flow] (update) -- node[lab,left] {下一轮参数} (policy);
\end{tikzpicture}
